\SourceLeaf{643}
\section*{REMARKS.}
\addcontentsline{toc}{section}{Remarks}
\par\nopagebreak[4]\smallskip\noindent\makebox[\linewidth][c]{\rule{.08\linewidth}{.35pt}}\par

The investigations collected here in 21 numbered items were recorded in fragments and date from periods quite far apart. Nos. 1 and 2 occur in a diary among records of observations made in March, June and July 1833 on the galvanic currents induced by magnets. Nos. 6, 12, 18 and 21 are on separate sheets. They cannot be dated precisely, except for No. 18, which bears a date. The remaining items, excluding the letter [20.], are reproduced here in the order in which they appear in a notebook, where numerous developments on quite different subjects occur between them. The last of these items, containing the proof of \textsc{Ampère}'s fundamental theorem, was entered only after 1843. The others appear to belong to the period from 1833 to 1836.

All the different forms adopted here for the law of interaction between galvanic current elements follow from the principle that magnetism and galvanic currents are interchangeable, which is emphasized especially in No. 20. The investigations by \textsc{Wilhelm Weber} mentioned in this letter laid the groundwork for a theory set out in its completed form in 1846 in the first paper on electrodynamic measurements. In that theory, the total interaction between two particles of electricity (with the corresponding signs assigned to $e$, $e'$), at a distance $r$ from each other, is measured by
\[
ee'\left(\frac1{rr}-\frac1{cc}\frac1{rr}\frac{\mathrm dr^2}{\mathrm dt^2}+\frac2{cc}\frac1r\frac{\mathrm ddr}{\mathrm dt^2}\right)
\]
A positive value of this quantity signifies repulsion and a negative value signifies attraction. In the expression, $t$ denotes time and $c$ a speed. In their investigations (1855) into expressing measurements of current intensity in mechanical units, \textsc{Kohlrausch} and \textsc{Weber} found this speed to be $439450.10^6\,\dfrac{\text{millimeter}}{\text{second}}$.

\SourceLeaf{644}

The theorem in No. 4 has been given a purely geometric form. I felt I had to give it this place because of its importance for the theory of galvanic currents. Consider the integral that determines how many times the closed curve $s$ links the system of closed curves $s'$. If the elements of all the closed galvanic currents present in space are substituted for $\mathrm ds'$, this integral gives the algebraic sum of the intensities of those currents that pierce a surface bounded by $s$ and otherwise chosen arbitrarily. With this interpretation of $\mathrm ds'$, the integral itself equals $\dfrac1{4\pi}\int\dfrac{\mathrm dV}{\mathrm ds}\mathrm ds$ if $V$ denotes the potential function for the magnetic effects of the currents $s'$, as in No. 9. The theorem thus corresponds to the one stated by \textsc{Gauss} in his treatise on the attraction of ellipsoids. That theorem determines the mass enclosed by a closed surface $(\omega)$ from the inward normal forces of its attraction $\left(\dfrac{\mathrm dV}{\mathrm dN}\right)$ at the surface, by means of $\dfrac1{4\pi}\int\dfrac{\mathrm dV}{\mathrm dN}\mathrm d\omega$.

One way to obtain the indicated value of the integral above, $\dfrac1{4\pi}\allowbreak\int\dfrac{\mathrm dV}{\mathrm ds}\mathrm ds$, is to convert the integral expressions for the derivatives of $V$ with respect to the coordinates into integrals over surfaces bounded by the individual current conductors $s'$, with the surfaces $\omega$ chosen arbitrarily. A theorem indicated in Art. 38 of the general theory of terrestrial magnetism gives the difference between the values of the potential function at corresponding points on the two sides of a surface suitably covered with opposite magnetic fluids. By that theorem, the form thus obtained for the derivative of $V$ with respect to $s$ shows immediately that the required integral equals the algebraic sum of the intensities of the currents flowing along the boundaries of the surfaces that the curve $s$ pierces, namely the surfaces $\omega'$.

Integrals over a closed curve $s$ can be converted into integrals over a surface bounded by $s$, denoted by $\omega$. The following theorem gives this conversion. Take any orthogonal rectilinear or curvilinear coordinates $\xi,\eta,\zeta$, so that the square of the length element is given in general by an expression of the form
\[
\xi'\xi'\mathrm d\xi^2+\eta'\eta'\mathrm d\eta^2+\zeta'\zeta'\mathrm d\zeta^2
\]
and suppose that, at the points of the surface $\omega$, the arbitrary functions $\lambda,\mu,\nu$ of the coordinates $\xi,\eta,\zeta$ vary continuously together with their first derivatives. Then the identity
\[
\begin{aligned}
&\int\left(\lambda\frac{\mathrm d\xi}{\mathrm ds}+\mu\frac{\mathrm d\eta}{\mathrm ds}+\nu\frac{\mathrm d\zeta}{\mathrm ds}\right)\mathrm ds\\
&=\int\left\{\left(\frac{\mathrm d\mu}{\mathrm d\zeta}-\frac{\mathrm d\nu}{\mathrm d\eta}\right)\frac{\xi'}{\eta'\zeta'}\frac{\mathrm d\xi}{\mathrm dn}+\left(\frac{\mathrm d\nu}{\mathrm d\xi}-\frac{\mathrm d\lambda}{\mathrm d\zeta}\right)\frac{\eta'}{\zeta'\xi'}\frac{\mathrm d\eta}{\mathrm dn}\right.\\
&\qquad\left.+\left(\frac{\mathrm d\lambda}{\mathrm d\eta}-\frac{\mathrm d\mu}{\mathrm d\xi}\right)\frac{\zeta'}{\xi'\eta'}\frac{\mathrm d\zeta}{\mathrm dn}\right\}\mathrm d\omega
\end{aligned}
\]
always holds when $\xi,\eta,\zeta$ in the first integral are the coordinates of a point of the length element $\mathrm ds$, and $\xi,\eta,\zeta$ in the second are the coordinates of a point of the surface element $\mathrm d\omega$, with $n$ the normal to that surface element. The positive direction of the normal must be chosen as follows. Let $\mathrm dt$ be the first element of a curve that starts at a point of $\mathrm ds$, is normal to that element, and lies in the surface $\omega$. It must be possible, by continuously moving the coordinate system in space, to bring the positive directions of $\xi'\mathrm d\xi$, $\eta'\mathrm d\eta$, $\zeta'\mathrm d\zeta$ into coincidence, in that order, with the positive directions of $\mathrm dn$, $\mathrm ds$, $\mathrm dt$.

\SourceLeaf{645}

Repeated application of this theorem also proves \textsc{Ampère}'s fundamental theorem in the general form that directly reduces the potential function for the interaction between surfaces covered in a specified way with magnetic fluid to the potential function for the interaction between galvanic currents whose positions and intensities are determined by those surfaces and the magnetization.

The proof in No. 9 that the different expressions for the potential function $V$ have equal values can be given a symmetric form by introducing the function $R$
\[
\begin{aligned}
={}&xx\operatorname{arctang}\frac{yz}{xr}+yy\operatorname{arctang}\frac{zx}{yr}+zz\operatorname{arctang}\frac{xy}{zx}\\
&+2yzi\operatorname{arctang}\frac{xi}r+2zxi\operatorname{arctang}\frac{yi}r+2xyi\operatorname{arctang}\frac{zi}r
\end{aligned}
\]
where $i$ stands for $\BareRoot{-1}$, and noting that the equations
\[
\begin{array}{ll}
\dfrac{\mathrm ddR}{\mathrm dx^2}=2\operatorname{arctang}\dfrac{yz}{xr},&\dfrac{\mathrm ddR}{\mathrm dy\mathrm dz}=2i\operatorname{arctang}\dfrac{xi}r\\[5pt]
\dfrac{\mathrm ddR}{\mathrm dy^2}=2\operatorname{arctang}\dfrac{zx}{yr},&\dfrac{\mathrm ddR}{\mathrm dz\mathrm dx}=2i\operatorname{arctang}\dfrac{yi}r\\[5pt]
\dfrac{\mathrm ddR}{\mathrm dz^2}=2\operatorname{arctang}\dfrac{xy}{zr}.&\dfrac{\mathrm ddR}{\mathrm dx\mathrm dy}=2i\operatorname{arctang}\dfrac{zi}r
\end{array}
\]
\[
\frac{\mathrm ddR}{\mathrm dx^2}+\frac{\mathrm ddR}{\mathrm dy^2}+\frac{\mathrm ddR}{\mathrm dz^2}=(2m+1)\pi,\qquad\frac{\mathrm d^3R}{\mathrm dx\mathrm dy\mathrm dz}=-2\frac1r
\]
hold. The derivatives of the function $R(x,y,z)$ also give finite expressions for the potential functions for the magnetic effect of galvanic currents whose conductors consist of straight segments parallel to the axes of a rectangular coordinate system.

We owe our knowledge of the arrangement of length comparisons that \textsc{Gauss} used to determine an interval of a graduated scale in units of the whole scale to a note that Privy Court Councillor \textsc{Weber} made for himself in 1839 or 1840.

The theoretical investigations into diffraction phenomena were probably prompted by the work on that subject published by \textsc{F. M. Schwerd} in 1835. The two general formulas for the effect of a luminous point $P$ on a point $p$ are different. Consider surface integrals whose elements do not depend on the position of the plane passing through a point of the corresponding surface element $\mathrm ds$ and through the points $P$ and $p$. In the corresponding curve integrals, the elements likewise do not depend on the position of the plane through a point of the corresponding element of the boundary line $u$ and through the points $P$ and $p$. Their differentials, however, represent a change only in the angle $\theta$ between that plane and a fixed plane through $P$ and $p$. The general conversion of the surface integrals into the curve integrals follows from the equation
\SourceLeaf{646}
\[
\begin{aligned}
\int Q\mathrm d\theta=\int\frac{h.Q.\sin v.\mathrm du}{rR\sin w}
={}&\int\frac h{rR}.\frac1{2\sin\frac12w^2}.\frac{\mathrm dQ}{\mathrm d(R-r)}\frac{\mathrm d(R+r)}{\mathrm d\rho}.\mathrm ds\\
&-\int\frac h{rR}.\frac1{2\cos\frac12w^2}.\frac{\mathrm dQ}{\mathrm d(R+r)}.\frac{\mathrm d(R-r)}{\mathrm d\rho}\,\mathrm ds
\end{aligned}
\]
This equation is a special case of the theorem mentioned in the preceding remark, provided that $Q$ denotes a quantity depending only on $R$ and $r$ and varying continuously, together with its partial derivatives with respect to $R+r$ and $R-r$, at the points of the surface $s$. In addition, $w$ is defined more precisely than in the text. It denotes the angle between $R$ and $r$ enclosed by the positive directions of those lines, taken to agree with the directions in which the light ray travels.

\begin{flushright}\textsc{Schering.}\end{flushright}

\par\nopagebreak[4]\medskip\noindent\makebox[\linewidth][c]{\rule{.20\linewidth}{.35pt}}\par
\clearpage
